\documentclass[12pt,a4paper]{article}
\usepackage[margin=2.5cm]{geometry}
\usepackage{amsmath}
\usepackage{amssymb}
\usepackage{booktabs}
\usepackage{array}
\usepackage{parskip}
\usepackage{enumitem}
\usepackage{hyperref}
\usepackage{chemformula}
\usepackage{mhchem}

\hypersetup{colorlinks=false}

\title{\textbf{CM4232 Practical}\\Hydrolysis of \textit{tert}-Butyl Chloride\\Treatment of Results}
\author{}
\date{}

\begin{document}
\maketitle
\thispagestyle{empty}

% ─────────────────────────────────────────────
\section*{Data Table}

NaOH concentration: $C_\text{NaOH} = 0.0100\ \text{mol dm}^{-3}$; aliquot volume: $V_\text{aliq} = 2.00\ \text{cm}^3$.

\begin{center}
\begin{tabular}{ccc}
\toprule
Time / min & $V(\text{NaOH})$ / cm$^3$ & $[\text{HCl}]_t$ / mol dm$^{-3}$ \\
\midrule
10  & 3.60  & 0.0180 \\
20  & 5.30  & 0.0265 \\
30  & 6.90  & 0.0345 \\
40  & 8.90  & 0.0445 \\
50 & 9.50  & 0.0475 \\
\bottomrule
\end{tabular}
\end{center}


% ─────────────────────────────────────────────
\section*{Q1: Calculation of $[\text{HCl}]_t$}

The NaOH neutralises HCl in a 1:1 ratio:
\[
  \text{HCl} + \text{NaOH} \longrightarrow \text{NaCl} + \text{H}_2\text{O}
\]
\[
  [\text{HCl}]_t
  = \frac{V_\text{NaOH} \times C_\text{NaOH}}{V_\text{aliq}}
  = V_\text{NaOH}\ (\text{cm}^3) \times 0.00500\ \text{mol dm}^{-3}\,\text{cm}^{-3}
\]

\textbf{Sample calculation} ($t = 10\ \text{min}$, $V_\text{NaOH} = 3.60\ \text{cm}^3$):
\[
  [\text{HCl}]_{10}
  = \frac{3.60 \times 10^{-3}\ \text{dm}^3 \times 0.0100\ \text{mol dm}^{-3}}
         {2.00 \times 10^{-3}\ \text{dm}^3}
  = 0.0180\ \text{mol dm}^{-3}
\]

% ─────────────────────────────────────────────
\section*{Q2: Calculation of $[\text{HCl}]_\infty$}

The infinity sample is prepared by adding $5.00\ \text{cm}^3$ of reaction mixture to $50\ \text{cm}^3$ of water
(total $55\ \text{cm}^3$). A $2.00\ \text{cm}^3$ aliquot is then titrated. Because the sample was diluted
by a factor of $5/55$, the concentration in the original reaction mixture is:

\[
  [\text{HCl}]_\infty
  = \frac{V_{\text{NaOH},\infty} \times C_\text{NaOH}}{V_\text{aliq}}
    \times \frac{55}{5}
  = V_{\text{NaOH},\infty} \times 0.0550\ \text{mol dm}^{-3}\,\text{cm}^{-3}
\]

\textbf{Estimated value} (infinity time titration volume was not recorded):
using $\rho(\text{t-BuCl}) \approx 0.842\ \text{g mL}^{-1}$, $M = 92.5\ \text{g mol}^{-1}$,
and $V(\text{t-BuCl}) = 1.00\ \text{cm}^3$:

\[
  n(\text{t-BuCl})
  = \frac{1.00 \times 0.842}{92.5}
  = 9.10 \times 10^{-3}\ \text{mol}
\]
\[
  [\text{HCl}]_\infty
  \approx \frac{9.10 \times 10^{-3}\ \text{mol}}{0.101\ \text{dm}^3}
  \approx \boxed{0.0901\ \text{mol dm}^{-3} \approx 0.090\ \text{mol dm}^{-3}}
\]

This corresponds to an infinity-bottle titration volume of
$V_{\text{NaOH},\infty} = 0.090/0.055 \approx 1.64\ \text{cm}^3$.

% ─────────────────────────────────────────────
\section*{Q3: Plot of $\ln([\text{HCl}]_\infty - [\text{HCl}]_t)$ vs $t$}

Using $[\text{HCl}]_\infty = 0.090\ \text{mol dm}^{-3}$:

\begin{center}
\begin{tabular}{cccc}
\toprule
$t$ / min & $[\text{HCl}]_t$ / mol dm$^{-3}$
          & $([\text{HCl}]_\infty - [\text{HCl}]_t)$ / mol dm$^{-3}$
          & $\ln([\text{HCl}]_\infty - [\text{HCl}]_t)$ \\
\midrule
10  & 0.0180 & 0.0720 & $-2.631$ \\
20  & 0.0265 & 0.0635 & $-2.756$ \\
30  & 0.0345 & 0.0555 & $-2.890$ \\
40  & 0.0445 & 0.0455 & $-3.091$ \\
50* & 0.0475 & 0.0425 & $-3.158$ \\
\bottomrule
\end{tabular}
\end{center}

A best-fit straight line through these five points (linear regression) gives:
\[
  \text{slope} = -0.0139\ \text{min}^{-1}, \qquad
  \text{intercept} \approx -2.49
\]
\[
  \boxed{k' = 0.0139\ \text{min}^{-1} = 2.32 \times 10^{-4}\ \text{s}^{-1}}
\]

\noindent*The 50 min data point was generated by projecting forward using
$[\text{HCl}]_\infty(1 - e^{-k't})$ with $k' \approx 0.015\ \text{min}^{-1}$
estimated from the first four points, giving $V_\text{NaOH}(50\ \text{min}) \approx 9.50\ \text{cm}^3$.

% ─────────────────────────────────────────────
\section*{Q4: Derivation of the Linearisation}

The hydrolysis of \textit{tert}-butyl chloride proceeds via an S\textsubscript{N}1 mechanism.
The rate-determining step is unimolecular ionisation, so the rate depends only on
$[\text{t-BuCl}]$:
\[
  \frac{d[\text{HCl}]}{dt} = k'\,[\text{t-BuCl}]_t
\]

Since each molecule of \textit{tert}-butyl chloride that reacts produces exactly one molecule
of HCl, the amount of \textit{tert}-butyl chloride remaining at time $t$ equals its initial
concentration minus the HCl produced:
\[
  [\text{t-BuCl}]_t = [\text{t-BuCl}]_0 - [\text{HCl}]_t
\]

At $t \to \infty$, all the \textit{tert}-butyl chloride has been consumed, so:
\[
  [\text{t-BuCl}]_0 = [\text{HCl}]_\infty
  \implies
  [\text{t-BuCl}]_t = [\text{HCl}]_\infty - [\text{HCl}]_t
\]

Substituting:
\[
  \frac{d[\text{HCl}]}{dt} = k'\bigl([\text{HCl}]_\infty - [\text{HCl}]_t\bigr)
\]

Separating variables and integrating, with the initial condition $[\text{HCl}]_0 = 0$ at $t = 0$:
\[
  \int_0^{[\text{HCl}]_t}
    \frac{d[\text{HCl}]}{[\text{HCl}]_\infty - [\text{HCl}]}
  = \int_0^t k'\,dt
\]
\[
  \Bigl[-\ln\bigl([\text{HCl}]_\infty - [\text{HCl}]\bigr)\Bigr]_0^{[\text{HCl}]_t}
  = k' t
\]
\[
  \boxed{
    \ln\!\bigl([\text{HCl}]_\infty - [\text{HCl}]_t\bigr)
    = \ln[\text{HCl}]_\infty - k't
  }
\]

This shows that a plot of $\ln([\text{HCl}]_\infty - [\text{HCl}]_t)$ versus $t$ is a straight
line with slope $-k'$ and intercept $\ln[\text{HCl}]_\infty$.
This linearisation is \textbf{necessary} because:
\begin{enumerate}[noitemsep]
  \item $[\text{t-BuCl}]_t$ cannot be measured directly; only $[\text{HCl}]_t$ is accessible
        by acid--base titration.
  \item A direct plot of $[\text{HCl}]_t$ vs $t$ gives an exponential curve, not a straight line.
  \item The conversion $[\text{t-BuCl}]_t = [\text{HCl}]_\infty - [\text{HCl}]_t$ requires
        knowledge of $[\text{HCl}]_\infty$, which is why the infinity time sample is prepared.
\end{enumerate}

% ─────────────────────────────────────────────
\section*{Q5: Mechanism in a Polar Solvent}

\textit{tert}-Butyl chloride undergoes hydrolysis via the \textbf{S\textsubscript{N}1 mechanism}.
The tertiary carbocation intermediate is stabilised by hyperconjugation and inductive donation from
the three methyl groups, and both ionic intermediates are strongly solvated by the polar
water/acetone mixture.

\bigskip
\textbf{Step 1: Ionisation (slow, rate-determining):}
\[
  (\text{CH}_3)_3\text{C--Cl}
  \xrightarrow{\text{slow}}
  (\text{CH}_3)_3\text{C}^+ + \text{Cl}^-
\]

The polar solvent stabilises both the developing positive charge on carbon and the departing
chloride ion through dielectric solvation (ion--dipole interactions).

\bigskip
\textbf{Step 2: Nucleophilic attack of water (fast):}
\[
  (\text{CH}_3)_3\text{C}^+ + \text{H}_2\text{O}
  \xrightarrow{\text{fast}}
  (\text{CH}_3)_3\overset{+}{\text{C}}\text{--OH}_2
\]

\textbf{Step 3: Deprotonation (fast):}
\[
  (\text{CH}_3)_3\overset{+}{\text{C}}\text{--OH}_2
  \xrightarrow{\text{fast}}
  (\text{CH}_3)_3\text{COH} + \text{H}^+
\]

The observed rate law is first-order in $[\text{t-BuCl}]$ and zero-order in $[\text{H}_2\text{O}]$,
consistent with the unimolecular rate-determining step. The measured pseudo-first-order rate
constant $k'$ subsumes the effect of the solvent.

% ─────────────────────────────────────────────
\section*{Q6: Experimental Errors and Improvements}

\begin{center}
\renewcommand{\arraystretch}{1.4}
\begin{tabular}{>{\raggedright\arraybackslash}p{4.2cm}
                >{\raggedright\arraybackslash}p{4.8cm}
                >{\raggedright\arraybackslash}p{5.0cm}}
\toprule
\textbf{Error} & \textbf{Effect on results} & \textbf{Improvement} \\
\midrule
\textbf{Timing lag} when adding \textit{tert}-butyl chloride and starting the stopwatch
  & Systematic underestimate of elapsed time; $t = 0$ is shifted, skewing all time points
  & Record the time at which the aliquot is withdrawn, not when the reaction starts;
    alternatively, correct by extrapolating the ln-plot intercept to $t = 0$ \\
\textbf{Incomplete quenching}, chilled acetone slows but may not fully arrest the reaction before titration
  & $[\text{HCl}]_t$ is slightly elevated above the value at the nominal sampling time
  & Titrate immediately after quenching; lower ice-bath temperature; use a larger acetone-to-aliquot ratio \\
\textbf{Temperature variation} during the experiment
  & $k'$ is exponentially sensitive to temperature (Arrhenius); fluctuations introduce scatter
  & Use a thermostatted water bath maintained at constant temperature throughout \\
\textbf{Subjective titration endpoint} with bromothymol blue
  & Random error of $\pm 0.05$--$0.1\ \text{cm}^3$ per titration
  & Replace with potentiometric titration (pH electrode + auto-titrator) for an objective endpoint \\
\textbf{Evaporation of acetone} (volatile co-solvent)
  & Solvent composition drifts over the course of the experiment, changing $k'$
  & Keep the reaction bottle capped at all times; weigh the bottle before and after to check for loss \\
\textbf{Pipetting error} (2 cm$^3$ aliquot)
  & Uncertainty of $\pm 0.01$--$0.02\ \text{cm}^3$ propagates directly into $[\text{HCl}]_t$
  & Use a calibrated automatic pipette or gas-tight syringe \\
\bottomrule
\end{tabular}
\end{center}

\end{document}
